\section*{Introduction and main results}
A collision bit is an output, not a description of the entire
experiment. Here the command specifies a spatial probability density,
a direction and a duration. All attempted preparations remain in the
denominator, including an attempted start in the unknown solid. This
convention makes opposed launch laws measure a signed occupation
difference rather than a conditional free-phase probability. The exact
data are functionals of the spatial input; a finite experiment uses a
growing spatial grid. No claim about uniformly prepared count laws is
made by reducing the per-attempt output to a bit.

\begin{maintheorem}
On the bounded periodic smooth class $\mathcal B_\eta$ of
Definitions~\ref{def:class} and \ref{def:margin}, with a known positive
patch margin $\eta$, a finite scalar experiment and a finite geometric
algorithm recover a smooth periodic presentation to $C^2$ error
$O(\nu)$, its primitive orbit count and the exact rational relations
of its period generators, with confidence $1-\delta$.
Each attempt returns one bit: a free segment hit is one; a solid start
and a free miss are both zero, and all attempts remain in the denominator.
The design uses $O(\nu^{-3})$ localized spatial commands of scale
$\sigma\asymp\nu^{3/2}$ and
$O(\nu^{-3}\log(C/(\nu\delta)))$ attempts. In the coupling model of
Theorem~\ref{thm:calibration}, combined position, timing and angular
displacement error at most $c\sigma$ is allowed. The arithmetic bound
is $O(\nu^{-6})$ in the model of Theorem~\ref{thm:constructive}; the
real period basis is estimated, not recovered exactly. All constants
and admissible accuracies depend on the displayed priors and $\eta$.
\end{maintheorem}

The central fixed-displacement identities are
\[
 B_a(q)-B_{-a}(q+a)=\one_{\mathcal O}(q+a)-\one_{\mathcal O}(q),
 \qquad
 u(x)=-\min_{0\le k\le m}\sum_{j<k}\delta_a u(x+ja).
\]
The second identity assumes a zero of the nonnegative $u$ on each chain.
A separation bound and a component-diameter bound provide that zero,
including for smoothed occupation. Sections~\ref{sec:reversal}--\ref{sec:finite}
retain the complete fixed-displacement inverse and its period-recognition
and existence-based finite-observation consequences. The principal
extensions in this article are as follows.

\smallskip\noindent
\emph{Randomized preparation without exact collimation.}
For a known displacement distribution $\rho$, the translated opposed
protocol measures $g=(T_\rho-I)u$. If the associated walk hits
$\{u=0\}$ within $m$ steps, the recursion
\[
 V_0=0,\qquad V_{n+1}=\max(0,T_\rho V_n-g)
\]
recovers $u$ at step $m$, with uniform forcing-error amplification at
most $m$. Theorem~\ref{thm:stopped} supplies a geometric cone condition
ensuring this stopping bound and allows a fixed continuous spread of
nominal directions and durations. This exact extension requires the
specified coupled reverse preparation; it is not a claim that arbitrary
unknown jitter preserves the balance identity.

\smallskip\noindent
\emph{Quantified apparatus tolerance.}
For localized inputs of scale $\sigma$, Theorem~\ref{thm:calibration}
controls the bias under a coupled start-position error $\ell$, timing
error $\tau$ and angle error $\alpha$ by
$C(\ell+\tau+b\alpha)/\sigma$, when the nominal duration is at most
$b$ and the combined error is at most $\sigma$. The estimate includes
grazing directions and errors depending on the commanded start. It is
a geometric stability bound, not a calibration procedure for an apparatus.
In particular, scale-independent propagation of a mean error does not
remove the need to sharpen input calibration with $\sigma$.

\smallskip\noindent
\emph{A finite geometric reconstruction.}
On the bounded periodic class $\mathcal B_\eta$ of
Definitions~\ref{def:class} and \ref{def:margin}, with a \emph{known}
positive nonperiod patch margin $\eta$, Theorem~\ref{thm:constructive}
uses the same fixed-duration opposed scalar commands to construct
complete component approximations by thresholding, a finite proximity
graph and convex hulls. No search over a compact infinite-dimensional
shape class is used in this theorem. Compensated smoothing supplies
$C^2$ supports and the retained finite patch test locks the period
relations. The uniform error, probability and operation bounds depend
on $\eta$ and all stated geometric and smoothness priors. The operation
count is in the declared arithmetic and fixed-kernel evaluation model,
not a bit-complexity or optimality assertion. Exact rational relations
refer to an accepted generator basis whose real coordinates have small
error, not to exact recovery of arbitrary real numbers.

These conclusions leave the earlier existence regularization in
Theorem~\ref{thm:finite} unchanged and identify a separate constructive
route. They do not remove the bounded periodic-presentation hypothesis
or certify an unknown patch margin at a finite time. The exact local
inverse itself does not require periodicity. The closest observation
framework is hit probability and active geometric probing; the precise
capacity-functional relation and theorem-level comparisons are in
Section~\ref{sec:hit-comparison}.
